\section{Conclusion}
\label{sec:conclusion}
We reproduced Shaer, a fine-tuned LLM for controlled classical Arabic poetry, and examined what its meter metric
measures. Re-scoring the released poems reproduced every published number, and the untuned baseline reproduced
closely, but new generations from the released adapter and a full re-training scored about 15 points lower on
base-meter accuracy, which points to undocumented details of the generation pipeline. We then built a formal verifier
of Arabic prosody and showed that the published metric overstates validity: \val{mEzeroStrictValid}\% of the released
poems are strictly valid in meter, rhyme and length, against \val{mRefStrictValid}\% of the human poems for the same
prompts, \val{mEzeroNoReading}\% of the bayts that the classifier accepts contain a hemistich that cannot be read in the
requested meter at all, and validity collapses as poems get longer. These findings held across nested models,
alternative specifications, independent generation runs and control samples.

Two lessons reach beyond this system. Generators of formal verse should be evaluated, and rewarded, by instruments that
can reject, not only classify. And reproductions of LLM generation results need the generation precision, engine and
checkpoint, not only the weights and the training script. The verifier, the analysis code and all predictions are
public. Future work will complete the pre-registered extensions (hemistich search, self-training and the quality study),
obtain an expert-labelled gold set for the verifier, and carry the approach to Urdu and Persian verse, whose meters
derive from the same tradition.
